\begin{center}\begin{tikzpicture}[font=\small,>=stealth,line join=round]
\draw[gray!40,rounded corners=3pt] (0,0) rectangle (14.2,8.3);\node at (7.1,7.700000000000001) {例3：先把大数移回基本区间，再代入平方};
\begin{scope}[shift={(2.9000000000000004,2.033333333333333)},x=0.4cm,y=0.3833333333333333cm]\draw[->,gray] (-4.5,0)--(10.5,0) node[right]{$x$};\draw[->,gray] (0,-1)--(0,11) node[above]{$y$};\draw[blue,thick,solid,domain=-4:-3,samples=180] plot(\x,{(\x+6)^2});\draw[blue,thick,solid,domain=-3:3,samples=180] plot(\x,{(\x)^2});\draw[blue,thick,solid,domain=3:9,samples=180] plot(\x,{(\x-6)^2});\draw[blue,thick,solid,domain=9:10,samples=180] plot(\x,{(\x-12)^2});\fill[orange!85!black] (2,4) circle(2.2pt);\fill[orange!85!black] (8,4) circle(2.2pt);\fill[orange!85!black] (-3,9) circle(2.2pt);\fill[orange!85!black] (3,9) circle(2.2pt);\fill[orange!85!black] (9,9) circle(2.2pt);\end{scope}\node[anchor=west,align=left,text width=5.0cm] at (8.4,4.0) {基本段 $[-3,3)$ 上 $f=x^2$\\[9pt]按周期6延拓\\[9pt]$2024=337\times6+2$\\[9pt]$f(2024)=f(2)=4$\\[9pt]图中 $f(2)=f(8)=4$};
\node[text width=12.6cm,align=center] at (7.1,.65) {图中每段来自同一基本段的平移；接头处实际取值为9，两侧也趋近9。};
\end{tikzpicture}\end{center}
